Mô hình Polynomial cho thấy sự cải thiện so với OLS, với MSE giảm từ 0.5253
xuống 0.494 và $R^2$ tăng từ 0.6097 lên 0.633. Kết quả này cho thấy dữ liệu tồn
tại yếu tố phi tuyến và việc mở rộng không gian đặc trưng giúp mô hình biểu
diễn tốt hơn mối quan hệ giữa các biến.

Ngược lại, mô hình RBF không đạt được kết quả như kỳ vọng khi có MSE cao hơn
(0.6971) và $R^2$ thấp hơn (0.4821) so với OLS. Nguyên nhân có thể đến từ việc
lựa chọn siêu tham số (như số lượng tâm hoặc độ rộng kernel) chưa phù hợp, hoặc
bản chất dữ liệu không phù hợp với dạng biểu diễn của RBF.

Cuối cùng, mô hình Fourier thể hiện hiệu năng vượt trội so với tất cả các mô
hình còn lại. Cụ thể, mô hình này đạt MSE thấp nhất (0.3873), RMSE thấp nhất
(0.6223), MAE thấp nhất (0.4466) và $R^2$ cao nhất (0.7123). Điều này cho thấy
dữ liệu có thể chứa các thành phần mang tính chu kỳ hoặc lặp lại, và Fourier là
phương pháp phù hợp để khai thác đặc điểm này, thậm chí hiệu quả hơn cả
Polynomial trong trường hợp này.
